\section{Tipo de investigación}

\subsection{De acuerdo con la orientación o finalidad}

La investigación es \textbf{aplicada}, porque utiliza técnicas de series temporales, aprendizaje automático y desarrollo de software para atender un problema concreto: la falta de información anticipada y contextualizada para evaluar decisiones de siembra de palta Hass en La Libertad. El producto será un modelo predictivo integrado con módulos de rentabilidad, escenarios y análisis histórico.

\subsection{De acuerdo con la técnica de contrastación}

La investigación es \textbf{explicativa}, debido a que evaluará el efecto del uso de la solución sobre indicadores de apoyo a la decisión, y \textbf{predictiva} en su componente tecnológico, porque estimará precios FOB a seis semanas. La contrastación técnica comparará Random Forest, HistGradientBoosting, ElasticNet y Prophet mediante MAE, RMSE, MAPE y $R^2$.
